%% ---------------------------------------------------------------------------
%% Machine-Checkable Anatomy -- IEEE conference paper, v2 architecture
%%
%% Build:
%%     python paper/collect_v2.py --with-benchmark   -> generated/v2.tex
%%     pdflatex main_v2 && bibtex main_v2 && pdflatex main_v2 && pdflatex main_v2
%%
%% NOTHING IN THE RESULTS SECTION IS TYPED BY HAND. Every number arrives via
%% %% generated by paper/collect_v2.py -- do not edit by hand

\newcommand{\GatesSettlingRatio}{0.146}
\newcommand{\GatesPermutationGap}{0.00e+00}
\newcommand{\GatesGradientCoverage}{100.0}
\newcommand{\GatesSparsityAchieved}{5.1}
\newcommand{\GatesDgSparsity}{1.95}
\newcommand{\GatesAchLow}{79.4}
\newcommand{\GatesAchHigh}{98.4}
\newcommand{\GatesReplayGain}{14.5}
\newcommand{\GatesCa3Before}{0.6407}
\newcommand{\GatesCa3After}{0.9448}
\newcommand{\GatesWeakestAblation}{34.0}
\newcommand{\GatesGatesPassed}{113}
\newcommand{\GatesGatesFailed}{0}
\newcommand{\CheckpointEpochs}{25}
\newcommand{\CheckpointExamples}{867200}
\newcommand{\CheckpointMinutes}{119.5}
\newcommand{\ArchitectureTrainable}{9365549}
\newcommand{\ArchitectureFrozen}{75940116}
\newcommand{\ArchitectureTracts}{26}
\newcommand{\ArchitectureTimesteps}{12}
\newcommand{\ArchitectureLatent}{256}
\newcommand{\ArchitectureSparsity}{0.05}
, written by collect_v2.py from the actual gates and
%% the actual checkpoint. The paper's central argument is that brain-inspired
%% architecture claims should be machine-checked; a paper making that argument
%% with hand-transcribed numbers would refute itself.
%% ---------------------------------------------------------------------------

\documentclass[conference]{IEEEtran}
\IEEEoverridecommandlockouts

\usepackage{cite}
\usepackage{amsmath,amssymb}
\usepackage{graphicx}
\usepackage{booktabs}
\usepackage{xcolor}
\usepackage{url}

\newcommand{\todo}[1]{\textcolor{red}{[TODO: #1]}}
%% generated by paper/collect_v2.py -- do not edit by hand

\newcommand{\GatesSettlingRatio}{0.146}
\newcommand{\GatesPermutationGap}{0.00e+00}
\newcommand{\GatesGradientCoverage}{100.0}
\newcommand{\GatesSparsityAchieved}{5.1}
\newcommand{\GatesDgSparsity}{1.95}
\newcommand{\GatesAchLow}{79.4}
\newcommand{\GatesAchHigh}{98.4}
\newcommand{\GatesReplayGain}{14.5}
\newcommand{\GatesCa3Before}{0.6407}
\newcommand{\GatesCa3After}{0.9448}
\newcommand{\GatesWeakestAblation}{34.0}
\newcommand{\GatesGatesPassed}{113}
\newcommand{\GatesGatesFailed}{0}
\newcommand{\CheckpointEpochs}{25}
\newcommand{\CheckpointExamples}{867200}
\newcommand{\CheckpointMinutes}{119.5}
\newcommand{\ArchitectureTrainable}{9365549}
\newcommand{\ArchitectureFrozen}{75940116}
\newcommand{\ArchitectureTracts}{26}
\newcommand{\ArchitectureTimesteps}{12}
\newcommand{\ArchitectureLatent}{256}
\newcommand{\ArchitectureSparsity}{0.05}


\begin{document}

\title{Machine-Checkable Anatomy: Making Brain-Inspired\\
Architecture Claims Falsifiable}

\author{\IEEEauthorblockN{Rishiraj Chakraborty}}

\maketitle

\begin{abstract}
Brain-inspired architectures are usually validated on task accuracy, which
cannot distinguish a system whose anatomy does work from one where the
anatomical language is decorative. We present a multimodal architecture of four
functionally distinct cortical lobes, a thalamus with reticular gating, basal
ganglia, limbic structures and a cerebellar forward model, in which every
anatomical claim is expressed as an executable test. Parallelism is established
by permutation invariance rather than asserted: all $24$ orderings of the four
lobes produce bit-identical output. Convergence, sparsity, Dale's principle,
corticothalamic balance, pattern completion and hippocampal replay are each
measured rather than described, and every one of \ArchitectureTracts{} named
white-matter tracts is shown to change the output when severed. We report
\GatesPassed{} passing checks. Against a parameter-matched monolithic MLP the
architecture loses by 1.0 points on a 20-class text task and 6.0 points on
100-class images, at roughly a hundred times the compute, and we report that
plainly:
every module mattering and the whole being better are separate claims, and only
the first is established here. We also report what the discipline exposed: three
distinct failures that produced plausible accuracy while a module was
functionally dead, each invisible to task metrics and each caught only by a
structural test. We argue that this class of failure is the normal case in
brain-inspired work, and that accuracy alone cannot detect it.
\end{abstract}

\begin{IEEEkeywords}
neuromorphic architecture, modular networks, verification, cortical models,
ablation analysis
\end{IEEEkeywords}

\section{Introduction}

A recurring pattern in brain-inspired machine learning is that a system is
described in anatomical terms---lobes, thalamic relays, basal ganglia
gating---and then evaluated on a benchmark. If accuracy is competitive, the
anatomy is taken to be doing something. This inference does not hold. A network
can contain a component named ``basal ganglia'' that receives input, produces
output, trains without error, and contributes nothing whatsoever to the result,
because the rest of the system has routed around it. Accuracy cannot see this.
Neither can a loss curve.

We take the position that an anatomical claim is only meaningful if it is
falsifiable, and that the falsification should be executable. If a system claims
its lobes run in parallel, there should be a test that fails when they do not.
If it claims a tract carries information, cutting the tract should change the
output, and the amount should be reported.

This paper describes such a system and, equally, what building it under that
constraint revealed.

\subsection{Contributions}

\begin{enumerate}
\item An architecture of four inductively distinct cortical lobes settling
      recurrently over $T=\ArchitectureTimesteps{} timesteps with explicit
      conduction delays, plus thalamic, basal ganglionic, limbic, cerebellar and
      neuromodulatory subsystems. Systems containing these structures exist
      \cite{marti4,braincog}; what is new here is that each commitment is
      stated so that it can fail.
\item A gate suite in which every anatomical claim is an executable test,
      with \GatesPassed{} checks passing. Parallelism is proved by permutation
      invariance, conduction delay by when a severed pathway is first felt,
      sparsity and Dale's principle by measurement rather than by construction
      alone.
\item The observation that ablation measured as state displacement and
      ablation measured as accuracy come apart entirely -- $134\%$ against
      $0.0$ points for the same severed pathway. Ablation is well established
      for establishing necessity \cite{meyes2019}; which of these two
      quantities it reports decides the conclusion.
\item Four documented failure modes in which a module was inert or being
      corrupted while task accuracy and loss curves looked healthy. This is
      the paper's most transferable result, and one of them inverted our own
      headline finding.
\end{enumerate}

\section{Related Work}

Three lines of work bear directly on this paper, and the position we occupy is
defined as much by what they already do well as by what none of them do.

\subsection{Brain-inspired modular architectures}

Decomposing perception into specialised interacting modules on neuroscientific
grounds is an active line of work. Recent cortically-inspired architectures
operationalise modular specialisation, cross-modal integration and predictive
feedback, and demonstrate that cortical organisation principles improve
modularity in practice \cite{cortical2026}. That work is explicit that it
implements functional \emph{principles} rather than anatomy: it has no
thalamus, no basal ganglia and no conduction delays, and it states the absence
of explicit temporal dynamics as a limitation.

Systems that do implement named subcortical structures exist. MARTI-4 couples
neocortex, basal ganglia and thalamus with the thalamus as an information hub
\cite{marti4}, and BrainCog provides a spiking substrate for brain-inspired
cognitive engines \cite{braincog}. Broader surveys map brain regions onto AI
paradigms and argue for the correspondence \cite{surveyjin,agentagi}.

Our architecture belongs to this family and is not novel for containing these
structures. What differs is that every anatomical commitment is made
falsifiable: conduction delays are enforced rather than described, so
parallelism is provable by permutation; every projection is a named object, so
a single pathway can be severed; and the sparsity, sign constraints and
corticothalamic balance are measured rather than asserted.

\subsection{Ablation as verification}

Ablation is long established for establishing that a component is necessary,
and has been studied systematically in artificial networks
\cite{meyes2019}. The standard question it answers is whether removing a part
degrades \emph{performance}.

We ask a different one, and the difference produced this paper's most useful
result. Our sweep measures how far removing a part moves the system's internal
state, and we found that this comes apart from performance completely: severing
the arcuate fasciculus displaces the settled cortical state by $134\%$ and
costs $0.0$ accuracy points. Both numbers are correct. A part can be
load-bearing for the representation and irrelevant to the answer, and an
ablation protocol that reports only one of those will conclude the opposite of
an ablation protocol that reports the other.

\subsection{What we claim}

Not that anatomical structure improves accuracy -- our own benchmark says it
does not, on these tasks. The claim is narrower and, we think, more useful: a
brain-inspired architecture should state its commitments as executable tests,
because the alternative is a literature of systems whose anatomical vocabulary
cannot be checked. Section~V reports four failures found this way, each of
which produced entirely normal training behaviour while a component was inert
or being corrupted.

\section{Architecture}

\subsection{No forward pass}

The system has no pipeline order. All four lobes update simultaneously from
delayed copies of one another:
\begin{equation}
s_\ell[t] = (1-\alpha)\, s_\ell[t-1] + \alpha \, f_\ell\!\left(
\textstyle\sum_{\tau \in \mathrm{in}(\ell)} W_\tau\, s_{\mathrm{src}(\tau)}[t - d_\tau]\right)
\end{equation}
where every $d_\tau \ge 1$. No module may read the present state of any other,
which makes update order irrelevant by construction rather than by convention.
The leak $\alpha$ is a membrane time constant; without it the state is a fresh
function of the inputs at every step and never converges.

\subsection{Four operators, one input}

The claim that the lobes are four distinct networks rather than one network
instantiated four times is load-bearing, and it is only meaningful if each
reads the same input differently. Every sensory sequence is therefore delivered
to all four lobes, and each applies its own operator to it: occipital filters
locally with grouped convolution over real spatial positions, parietal attends
between positions, temporal carries order through recurrence, and frontal
scores positions for relevance and integrates by that score. The modalities
remain distinguished by the thalamic nucleus they enter and by which lobe is
listed first, not by which lobes are permitted to see them.

This was not the original arrangement, and the difference is measurable. When
each modality was routed to a single lobe, a text input reached only temporal
and the remaining three lobes carried nothing but echoes returning over
cortico-cortical tracts; three of the four networks were not parsing the input
at all. That is invisible in accuracy, which is the recurring theme of
Section~V.

\subsection{Named tracts}

Every connection is a first-class object with a source, a target, a conduction
delay and a lesion switch. There is no anonymous tensor addition anywhere in the
update. This is what makes single-pathway ablation possible, and it is a
constraint on the implementation rather than an annotation of it.

\subsection{Laminar organisation}

Afferents are routed by cortical layer: thalamic drive to L4, lateral and
feedforward cortico-cortical input to L2/3, feedback to L5/6, with L6 supplying
the corticothalamic projection. Each tract declares its layer, and a tract
without a declared layer raises an error at construction. Superficial layers
enforce $k$-winners-take-all sparsity at
\ArchitectureSparsity{} and separate excitatory and inhibitory populations whose
weights are held non-negative by construction, so Dale's principle is a property
of the parameterisation rather than a hope about training.

\subsection{Subcortical systems}

The basal ganglia implement direct, indirect and hyperdirect pathways with a
tonically active output nucleus, so ``Go'' operates by disinhibition. Critically,
their output goes to the thalamus and returns to frontal cortex; it does not
gate the output. Dopamine from the substantia nigra scales the gradient reaching
striatal weights rather than scaling activations, making it a learning rate
rather than a gain. The limbic system runs on two further timescales: an
amygdala with fast and slow routes separated by conduction delay, a hippocampus
implementing separation, completion and comparison as distinct stages, and a
hypothalamus whose state persists across episodes. A cerebellar forward model
predicts the next cortical state, with an inferior olive carrying only the
difference.

\section{Verification Methodology}

Each claim in Section II is paired with a test that can fail. We summarise the
non-obvious ones.

\subsection{Parallelism by permutation}

Rather than asserting that the lobes are parallel, we execute all $4! = 24$
orderings of the lobe update and require bit-identical output. The measured
maximum difference is \GatesPermutationGap{}. A delay of zero is separately
rejected at construction, since it would reintroduce order dependence.

\subsection{Convergence, not termination}

A recurrent system that stops at $T$ has not necessarily settled. We measure the
change in cortical state per timestep and require both a downward trend and a
final change bounded relative to the first, with a leak of $1.0$ retained as a
control that must fail. The measured settling ratio is \GatesSettlingRatio{}.

\subsection{Gradient coverage with declared exceptions}

We require every trainable tensor to receive gradient, measured across a full
epoch rather than a single batch---a single-batch sample reports every pathway
not exercised by that batch as dead. Coverage is
\GatesGradientCoverage{}\%. Tensors that are structurally untrainable are listed
by name with the reason, so an \emph{undeclared} gap fails the gate.

\subsection{Ablation with named exceptions}

Every tract, subsystem and lobe is removed individually and the resulting shift
in settled state is reported. The weakest live part moves the output by
\GatesWeakestAblation{}\%. Two parts are quiet at rest and are declared: the
hyperdirect brake, measured while nothing requires braking, and the
neuromodulatory dial, which rests at an exactly neutral centre by construction.
Both have separate gates demonstrating function when driven.

\section{Results}

\subsection{Verified properties}

\begin{table}[t]
\caption{Measured architectural properties}
\label{tab:properties}
\centering
\begin{tabular}{lr}
\toprule
Property & Measured \\
\midrule
Lobe update orderings, bit-identical & 24 / 24 \\
Settling ratio (final / first change) & \GatesSettlingRatio{} \\
Cortical sparsity achieved & \GatesSparsityAchieved{}\% \\
Dentate gyrus sparsity achieved & \GatesDgSparsity{}\% \\
Gradient coverage & \GatesGradientCoverage{}\% \\
Feedforward share, low ACh & \GatesAchLow{}\% \\
Feedforward share, high ACh & \GatesAchHigh{}\% \\
CA3 completion (cue $\rightarrow$ recalled) & \GatesCaThreeBefore{} $\rightarrow$ \GatesCaThreeAfter{} \\
Replay retention gain & $+$\GatesReplayGain{} pts \\
Weakest live ablation & \GatesWeakestAblation{}\% \\
\bottomrule
\end{tabular}
\end{table}

Table~\ref{tab:properties} reports the measured values. Of these, the
acetylcholine result is the clearest demonstration that a neuromodulator is
doing what it is named for: a single scalar moves the feedforward share of
cortical drive from \GatesAchLow{}\% to \GatesAchHigh{}\%, monotonically, while
being exactly neutral at its baseline.

\subsection{Comparison against a monolithic baseline}

We compare against a monolithic MLP holding $98.6\%$ of the modular system's
trainable parameters, on identical frozen features, with the same optimiser and
the same seeds. Crucially, both are trained to CONVERGENCE rather than for equal
epochs: scored at six epochs each the comparison read $41.2\%$ against $2.9\%$,
which measures convergence speed and nothing else. Given a budget in which both
plateau, and scoring each on its best held-out epoch:

\begin{table}[h]
\caption{Modular architecture against a parameter-matched MLP.
Two seeds, 10000 examples per task, best held-out epoch, identical inputs and identical cosine schedule.}
\label{tab:benchmark}
\centering
\begin{tabular}{lrrr}
\toprule
Task & Chance & Monolith & Modular \\
\midrule
CIFAR-100 (image) & 1.0\% & \textbf{59.4\%} & 53.4\% \\
20 Newsgroups (text) & 5.0\% & \textbf{65.0\%} & 64.0\% \\
\midrule
Wall clock per seed & & 49 s & 5117 s \\
\bottomrule
\end{tabular}
\end{table}

\textbf{The modular architecture loses, by $6.0$ points.} On text the
difference is smaller than the seed-to-seed spread and we decline to call it a
result in either direction. The modular system also costs roughly eighty-five
times more compute per run, which is the cost that does not shrink.

The deficit does shrink, and the trend is more informative than any single
row. Measured at two scales after the defect described below was fixed:

\begin{center}
\begin{tabular}{lrrr}
\toprule
Examples per task & Monolith & Modular & Gap \\
\midrule
2000 (with the defect) & 46.1\% & 21.8\% & 24.3 \\
2000 & 44.5\% & 39.1\% & 5.5 \\
5000 & 50.5\% & 47.3\% & 3.2 \\
10000 & 59.4\% & 53.4\% & 6.0 \\
\bottomrule
\end{tabular}
\end{center}

Both architectures improve substantially with data. We previously read the
first two rows as a shrinking deficit and said so; the third row falsifies
that reading, and we record the correction rather than the tidier trend. A
gap inferred from two points is not a trend, and the monolith gains from
scale at least as fast as the modular system does.

One asymmetry is worth stating because it cuts against us. The full training
pipeline optimises five objectives and reaches $59.1\%$ on images at 20000
examples; the benchmark trains both architectures on the task loss alone, so
the modular system is measured here without the auxiliary objectives it was
designed around. That is the fair comparison -- the monolith has no such
objectives to give it -- but it means the benchmark understates what the
complete system does, and does not rescue the comparison.

This table replaces an earlier one reporting a $24.3$-point deficit, and the
reason belongs in the paper rather than in a changelog. Most of that gap was a
bug. \texttt{nn.Module.train()} recurses into every child, so the frozen
ResNet core's ten BatchNorm layers were switched into training mode by every
call, updating their running statistics from cortical state instead of the
ImageNet activations those statistics describe. Freezing parameters with
\texttt{requires\_grad\_(False)} stops gradients and does nothing about
buffers. Fixed, image accuracy on the full training set moved from $1.8\%$ to
$48.1\%$, and the benchmark gap fell from $24.3$ points to $5.5$.

We report this prominently because it is this paper's thesis turned on the
paper itself. The corrupted runs produced entirely normal loss curves and no
tracked metric indicated a problem. The bug was found only by comparing against
an \emph{untrained} baseline and noticing that training made the vision
pathway worse than leaving it alone. Had the first table been published, a bug
would have been published as a finding about architecture.

We report this without qualification because it is the finding. It also sharpens
rather than undermines the paper's argument: every one of the
\ArchitectureTracts{} tracts in this system demonstrably changes the output when
cut, no module is decorative, and the architecture still loses to one MLP. Those
two facts are compatible, and holding them together is the point. ``Every part
matters'' and ``the whole is better'' are different claims, and a literature
that verifies neither routinely implies both.

The mechanism is not mysterious. Hard $k$-winners-take-all leaves roughly
thirteen active units per lobe, fifty-two across the four, against the
monolith's $1728$ dense units. For twenty topics that is sufficient and the two
architectures tie; for a hundred image classes it is not. The sparsity is a
verified property with its own gate, and we did not relax it to improve this
number.

\section{What Verification Exposed}

This section reports failures. We consider it the most useful part of the paper,
because each failure produced plausible training behaviour while a component was
inert, and none would have been caught by accuracy.

\subsection{A dead module behind healthy accuracy}

Bounding cortical state by dividing by its norm is a standard stabilisation.
It also makes the output scale-invariant, which projects away the gradient with
respect to magnitude. Several modules here act on magnitude rather than
direction---the basal ganglia scale their thalamic drive---so the striatal
gradient fell from $63$ to $3.6 \times 10^{-8}$. Training proceeded normally,
the loss fell, and held-out accuracy reached a respectable value with the
striatum entirely inert. Clamping only the excess magnitude preserves the bound
while leaving the gradient exact below the ceiling.

\subsection{Unbounded scale in a squared error}

Three separate objectives---predictive coding, the cerebellar forward model, and
the cerebellum's own reported error---computed squared error on cortical state
whose magnitude nothing constrained. Cross-entropy always reduces its loss by
enlarging logits, the cheapest route to which is inflating that state, after
which the squared errors grow quadratically. Measured, the forward objective
went $1.28 \rightarrow 1.61 \rightarrow 7.3\times10^{8}$ across three epochs.
Each instance looked like a different bug. Normalising each error by its
target's own energy removes the coupling.

\subsection{Sparsity as a capacity limit}

Hard $k$-winners-take-all at \ArchitectureSparsity{} leaves roughly thirteen
active units per lobe. On a 20-class task with $\sim$190 examples per class this
is ample; on a 100-class task with 40 examples per class the same read-out sat
at chance, while a linear probe on the identical frozen features reached
$56.8\%$. Reading all four lobes rather than one recovered part of the gap. We
report the remainder as the genuine cost of the sparsity constraint rather than
relaxing a verified property to improve a number.

\subsection{A frozen module that was not frozen}

Freezing a pretrained component is usually written as
\texttt{requires\_grad\_(False)} plus a call to \texttt{eval()}. That stops
gradients and stops nothing else. \texttt{nn.Module.train()} recurses into
every child, so each call on the enclosing model returned the frozen ResNet
core's BatchNorm layers to training mode, where they update running statistics
from whatever passes through -- cortical state, not the ImageNet activations
those statistics describe.

The failure is invisible and modality-specific. Loss curves are normal.
Vision degrades because BatchNorm carries running buffers; the transformer
cores are untouched because LayerNorm carries none. Measured, this was the
difference between $1.8\%$ and $48.1\%$ on CIFAR-100, and it survived several
architecture revisions because nothing being tracked could see it. It was
found by comparing against an untrained baseline: a linear probe on the
\emph{untrained} cortex scored $13.3\%$, which is the observation that
training was making the pathway worse than leaving it alone.

\subsection{Tests that mutate what they measure}

Two gates were found to be measuring their own side effects. A cortex left in
training mode commits an episode to the hippocampus on every settle, so a
pattern-completion test run afterwards measured interference rather than
completion, and a lesion-reversibility test attributed the resulting drift to an
irreversible lesion. Both now isolate their own state. We note this because a
verification suite is itself software and is not exempt from the standard it
enforces.

\section{Limitations}

The auditory and somatosensory pathways are trained on synthetic data
constructed for the purpose; accuracy on them measures nothing about the
architecture and is reported only to support the statement that those pathways
are trained rather than randomly initialised. The concept vocabulary available
to the language interface is drawn from the frozen encoders and extends beyond
the labels the read-out was trained against; retrieved concepts outside that set
are marked as extrapolation. Scale is modest throughout.

\section{Conclusion}

We have described a brain-inspired architecture in which anatomical claims are
executable tests, and reported \GatesPassed{} passing checks covering
parallelism, convergence, sparsity, Dale's principle, corticothalamic balance,
pattern separation and completion, replay, forward-model learning and
per-pathway ablation. On task accuracy it ties a parameter-matched MLP on text
and loses substantially on images, at far greater cost. We take that seriously:
the architecture is verified, not vindicated. The discipline's value as a method
and the architecture's value as a model are separate questions, and only the
first is settled here. The failures in Section V are the part we would most like to see
replicated: in each, a component was dead while the system looked healthy, and
only a structural test could tell.

\bibliographystyle{IEEEtran}
\bibliography{refs}

\end{document}
